\section{Discussion}

Metab: add something about methionine cycle here. does it phenocopy molecularly? Check lit. Obata and Miura (2015).
\eh{Move this here: "The convergence of transcriptomic and metabolomic datasets indicates an orchestrated remodeling of peripheral physiology imposed by the activation of a small number of R2 neurons in the brain."}

The central finding of this study is that direct manipulation of adult EB ring neurons is sufficient to extend lifespan in \textit{Drosophila}, supporting a model in which the EB acts instructively rather than merely as a passive conduit for lifespan-altering sensory information. Previous work showed that R2 and R4 neurons are required for the lifespan-shortening effects of death perception, establishing these cells as a neural substrate through which sensory experience can influence aging. Here, we show that direct activation of R2 neurons, co-activation of R2/R4 neurons, and inhibition of R4d neurons are each sufficient to promote longevity under standard laboratory conditions. Strikingly, these effects can be elicited by manipulating very small neuronal populations, in some cases fewer than ten neurons per hemisphere, underscoring the sensitivity of organismal aging to defined circuit states. Together, these findings argue that EB ring neurons are not simply necessary for transmitting environmental information relevant to aging, but can themselves generate signals that modulate lifespan.

A second major conclusion is that distinct EB manipulations converge on increased lifespan through divergent proximal molecular states. Bulk head transcriptomics after one week of treatment revealed that R2 activation, R4d inhibition, and R2/R4 co-activation do not produce a common downstream transcriptional program. Instead, R2 activation induced a starvation-like signature that included nutrient-responsive IIS components, autophagy-associated genes, and one-carbon and transsulfuration-associated pathways, whereas R2/R4 co-activation preferentially engaged biosynthetic and storage-associated programs, including amino acid metabolism, glycogen synthesis, and NADPH-generating pathways. R4d inhibition produced a more limited transcriptional response, but attenuated stress-associated changes observed in green-light controls. These results argue against a simple model in which all lifespan-extending EB manipulations act through the same stereotyped physiological state. Rather, they suggest that EB activity can engage multiple downstream programs, each compatible with enhanced longevity.

The transcriptional data also point to specific pathways that may link EB activity to aging-regulatory physiology. The induction of \textit{foxo}, \textit{Thor}, \textit{InR}, \textit{dilp6}, and \textit{Cbs} following R2 activation suggests engagement of nutrient-sensing, proteostatic, and sulfur-amino-acid metabolic processes previously implicated in lifespan control. At the same time, the R2-induced state does not resemble a canonical starvation program. Nutrient-responsive IIS components were induced alongside ribosome biogenesis and translational programs that are not ordinarily co-regulated in this way, suggesting that direct EB activation may uncouple pathways that are typically coordinated by nutrient availability. More broadly, these data are consistent with the possibility that EB activity influences aging by imposing non-canonical nutrient-state programs rather than by simply mimicking physiological starvation or satiety. Because prior work showed that R2/R4 neurons regulate \textit{dilp3} and \textit{dilp5} expression in insulin-producing cells during death perception, one plausible hypothesis is that EB-dependent internal states are translated into systemic aging effects through downstream neuroendocrine centers. However, the present data do not establish that mechanism directly, and defining the relevant effector pathways will require future work.

Notably, the molecular signatures induced by EB manipulation were not mirrored by simple changes in organismal behavior or bulk physiology. Feeding assays provided no evidence that the transcriptional states elicited by R2 or R2/R4 activation were secondary to altered food intake, arguing against dietary restriction or compensatory hunger as the basis of the observed longevity phenotypes. Likewise, changes in TAG stores did not predict starvation resistance, and reduced TAG could occur without increased locomotor activity. Negative geotaxis similarly failed to reveal early behavioral correlates of the transcriptional states measured after one week of manipulation. Together, these results indicate that EB-dependent regulation of aging is not readily explained by gross changes in feeding, energy expenditure, lipid depletion, or locomotor performance. Instead, they support a model in which EB activity imposes internal physiological states whose consequences are reflected more strongly in gene expression than in canonical behavioral or metabolic outputs measured under standard laboratory conditions.

The valence experiments add an important dimension to this model by showing that lifespan-extending EB states need not share the same immediate motivational sign. In the FLIC, R2 activation in females was aversive, promoting rapid disengagement from the sucrose well without suppressing re-approach, while R2/R4 co-activation was positively valenced, increasing interactions and return frequency without prolonging individual feeding bouts. In a feeding-independent place-preference assay, the positive valence of R2/R4 co-activation generalized across sex and stimulation regime, whereas the consequences of R2 or R4d manipulation were more context dependent. These findings indicate that EB ring neurons do not contribute fixed, additive aversive or appetitive signals. Rather, the motivational meaning of EB activity appears to depend on population combination, assay context, sex, and stimulation conditions. Particularly striking is the observation that both aversive and appetitive EB states can promote longevity. This argues against a simple model in which lifespan extension is coupled to one privileged valence state and instead suggests that distinct internal states may engage different downstream mechanisms that converge on slower aging.

This nonlinear relationship between EB population activity and behavioral output may be especially informative in light of known interactions among ring-neuron classes. Prior work suggests that R2, R4m, and R4d neurons participate in an incoherent feedforward circuit, raising the possibility that manipulations such as R2 activation, R2/R4 co-activation, and R4d inhibition may overlap in some network-level consequences even while producing distinct transcriptional and behavioral states. Our data do not resolve whether these interventions ultimately converge on a shared terminal mechanism or instead reach similar lifespan outcomes through partially distinct pathways. However, the marked divergence in both transcriptional state and valence strongly suggests that the relevant aging signal is not encoded by the activity of any single ring-neuron class in isolation. A more parsimonious interpretation is that EB circuits influence aging through population-level configurations whose downstream consequences depend on the particular combination of neurons engaged.

Our findings that discrete populations of ellipsoid body ring neurons can instructively modulate lifespan through the generation of distinct internal states carry potentially significant implications for understanding the neural basis of aging in humans. The \textit{Drosophila} central complex, of which the ellipsoid body is a core component, shares strong neuroanatomical and functional homology with the vertebrate basal ganglia. Both structures serve as integrative hubs that process sensory information, coordinate motor output, and—critically—contribute to the generation of motivational and affective states. Recent work has expanded our understanding of the basal ganglia beyond its classical role in movement control, implicating this structure in the recognition of emotional stimuli, inference of others' emotional states, and awareness of subjective well-being {Gendron, 2023 128}. In particular, limbic basal ganglia output pathways, including ventral pallidal circuits, can assign appetitive or aversive value to ongoing behavior and reshape responses to challenge in a state-dependent manner. Our findings raise the possibility that a similarly general principle applies to aging: circuit configurations that impose distinct motivational or physiological states may converge on long-term effects on organismal maintenance and survival. More generally, these results suggest that lifespan may be modulated not only by what animals perceive, but by how neural state transforms the meaning of those perceptions into coordinated physiological output. If analogous circuits in the human basal ganglia similarly encode valence states that influence systemic physiology, this would provide a mechanistic framework for understanding associations between psychological states, chronic stress, and aging in human populations.

%% AI wanted some limitations. lol.
Several limitations should be considered when interpreting these findings. First, our transcriptomic analyses were performed on female heads at a single early time point, and thus do not capture male-specific or later-emerging molecular states. This is especially relevant given the sex-specific behavioral effects observed in place preference and negative geotaxis. Second, bulk head RNA-seq cannot distinguish transcriptional changes arising within the brain from those occurring in peripheral head tissues, nor can it resolve cell-type-specific responses downstream of EB manipulation. Third, the green light required for GtACR1 complicated behavioral interpretation of some inhibition experiments, particularly in the FLIC, where abrupt light onset itself appeared to be startling. Finally, although the present work shows that EB manipulations are sufficient to alter lifespan and to impose distinct transcriptional and motivational states, it does not yet identify the downstream endocrine, metabolic, or peripheral effectors that execute these aging phenotypes. These limitations define clear priorities for future work, including cell-type-resolved transcriptional analysis, direct study of IPC and neuroendocrine outputs, and more refined dissection of sex-specific responses.